%% ma: L12-B08-0061
%% lop: 12
%% chuong: 2
%% bai: 8
%% dang: 12.8.36
%% loai: DS
%% muc: [NB, NB, TH, VD]
%% dap_an: "ĐĐSĐ"
%% nguon: "LaTeX giáo viên gửi 21/10/2026 (M05, Vectơ và hệ trục tọa độ trong không gian, 107 câu), Phần II Câu 17"
%% kien_thuc: "Bài 7 (hình chiếu lên trục, mặt phẳng tọa độ), Bài 8 (trung điểm, độ dài, tích vô hướng), lớp 10 (trực tâm tam giác vuông)"
%% trang_thai: chinh-thuc
%% ghi_chu: "Giáo viên đã duyệt 21/10/2026. Hình gốc vẽ sẵn E, F và các đường hình chiếu (gợi ý lời giải ý a), b), c)); bản kho chỉ vẽ hệ trục và tam giác ABC"
\begin{ex}
Trong không gian $Oxyz$, cho ba điểm $A(3;-2;5)$, $B(2;1;-3)$ và $C(-1;2;1)$ (tham khảo hình bên dưới).
\begin{center}
\begin{tikzpicture}[x={(-.7cm,-.45cm)},y={(.8cm,0cm)},z={(0cm,.8cm)},scale=.75]
  \draw[truc] (-3,0,0)--(4.5,0,0) node[below left]{$x$};
  \draw[truc] (0,-3,0)--(0,4,0) node[right]{$y$};
  \draw[truc] (0,0,-4)--(0,0,6.2) node[left]{$z$};
  \node[below right,font=\footnotesize] at (0,0,0) {$O$};
  \coordinate (A) at (3,-2,5); \coordinate (B) at (2,1,-3); \coordinate (C) at (-1,2,1);
  \draw (A)--(B)--(C)--cycle;
  \foreach \P in {A,B,C} \fill (\P) circle (1.4pt);
  \path (A) node[left]{$A$} (B) node[right]{$B$} (C) node[above right]{$C$};
\end{tikzpicture}
\end{center}
\choiceTF
{\True Hình chiếu vuông góc của điểm $A$ trên trục $Oz$ là $A'(0;0;5)$.}
{\True Hình chiếu vuông góc của điểm $B$ trên mặt phẳng $(Oxz)$ là $B'(2;0;-3)$.}
{Biết $E$ là điểm nằm trên trục $Ox$ và $F$ là điểm nằm trên mặt phẳng $(Oyz)$ sao cho $C$ là trung điểm của đoạn thẳng $EF$. Khi đó $EF>5$.}
{\True Gọi $I(a;b;c)$ là trực tâm của tam giác $ABC$. Khi đó $a+b+c=2$.}
\loigiai{a) $A'(0;0;5)$. Đúng.

b) $B'(2;0;-3)$. Đúng.

c) $E(e;0;0)$, $F(0;f;g)$, $C$ là trung điểm $EF$ nên $e=-2$, $f=4$, $g=2$. $EF=\sqrt{4+16+4}=\sqrt{24}\approx4{,}9<5$. Sai.

d) $\overrightarrow{CA}=(4;-4;4)$, $\overrightarrow{CB}=(3;-1;-4)$, $\overrightarrow{CA}\cdot\overrightarrow{CB}=12+4-16=0$ nên tam giác $ABC$ vuông tại $C$, trực tâm là $C$. $a+b+c=-1+2+1=2$. Đúng.}
\end{ex}
